\documentclass[11pt]{article}

\usepackage[a4paper,top=15mm,bottom=9mm,left=29mm,right=29mm]{geometry}
\usepackage{mathpazo}
\usepackage[T1]{fontenc}
\usepackage[utf8]{inputenc}
\usepackage{microtype}
\usepackage{amsmath,amssymb}
\usepackage{parskip}

\setlength{\parskip}{0.24em}
\linespread{1.0}

\pagestyle{empty}

\begin{document}

\begin{center}
  {\large\bfseries Vineyard monodromy in Newton's equations}\\[0.3em]
  {\small Shawn Ray}
\end{center}

\vspace{0.2em}

Looking for vineyard monodromy in the wild, I found it in classical mechanics,
in the figure-eight solution of the three-body problem. What it measures is the
more interesting half.

Three equal masses chase one another around a figure-eight curve, the orbit
Moore found numerically and Chenciner and Montgomery proved exists. Sum a radial
kernel at each body and follow the degree-zero persistence of its superlevel
sets over a third of the period, the essential class dying at the infimum. The
diagram returns to itself with its three points cyclically permuted, so the
monodromy has order $3$. Those three features stay clear of the diagonal
throughout, the least persistent at $0.06$, and are distinct except at the wall
crossings, where two coincide exactly; following the critical points through
shows them exchanged, not turned back, so the cycle is transport, not a
relabeling of coincident points. Ordering strands by birth at the first slice,
running time forwards and signing each crossing by the death coordinate, the
braid is $(\sigma_2^{-1}\sigma_1)^2$, whose trace closure is the figure-eight
knot, by its Alexander polynomial $t-3+t^{-1}$ and independently by SnapPy.

The loop is $T/3$ because after a third of a period the bodies occupy the same
three positions with their labels shifted, and the density is a \emph{sum} over
them, so relabeling leaves it unchanged as a function and the diagram returns
exactly, not nearly. Scanning for set-returns finds only $T/3$ and $2T/3$, so
$T/3$ is primitive; over the full period the relabeling is the identity and the
monodromy trivial. Nothing about the loop is chosen.

What orders the vines is easy to write down. At the bodies $h_i=\sum_j
K(r_{ij})$, the shared term cancels, and $h_1-h_2=K(r_{13})-K(r_{23})$ vanishes
when body 3 is equidistant from bodies 1 and 2, an \emph{isosceles
configuration}. The peaks do the same, and exactly. Let $\sigma$ be the
reflection in the perpendicular bisector of bodies 1 and 2. It fixes that pair,
so $\rho\circ\sigma-\rho=K(|x-\sigma q_3|)-K(|x-q_3|)$, which is positive on
body 2's side whenever body 3 lies nearer body 1. Then $h_1\ge\rho(\sigma
M_2)>\rho(M_2)=h_2$, the first step because $\sigma M_2$ sits as far from body 1
as $M_2$ does from body 2, and $\rho$ is concave on that ball, so $M_1$ is its
highest point; and at equidistance $\sigma$ carries $\rho$ to itself, so the two
are equal. So wherever that concavity holds, the walls are the isosceles
configurations and nothing else, and the strands are ordered by the opposite
sides. The deaths are the two shortest, the minimum spanning tree, measured
rather than proved. The braid reads only the order of the births and of the
deaths, never their size.

The condition contains no kernel, with a consequence I did not expect: a
Gaussian, a cusped exponential, a Cauchy and a Student $t$ all put the walls at
the same twelve instants. Nor can the bandwidth. What it sets is a window: below
$\sigma\approx0.24$ the crowding differences fall towards numerical resolution,
above $\sigma\approx0.30$ a feature meets the diagonal, neither bound
topological.

A change of frame, even a time-dependent one, carries $\rho(x,t)$ to
$\rho(g_t^{-1}x,t)$. Every superlevel set goes to a congruent one, so each
diagram is untouched and the whole vineyard with it. Viewed from six frames
rotating at different rates the vineyard braid is $(\sigma_2^{-1}\sigma_1)^2$
every time, while the trajectory braid picks up a twist per turn. The two are
invariants of different objects: the trajectory braid of the motion in the
plane, the vineyard braid of that motion reduced by translations and rotations.

So the monodromy is no coincidence of this orbit: its order is $3$ because the
time shift permutes the bodies cyclically, and over $T/3$ the curve crosses four
of the twelve walls, the braid being its itinerary through the chambers. And the
knot carries content of its own. Over that same $T/3$ the bodies trace the
pseudo-Anosov braid $\sigma_1\sigma_2^{-1}$, up to the conjugacy a projection
introduces; under the standard map to $SL(2,\mathbb{Z})$ it is
$\left(\begin{smallmatrix}2&1\\1&1\end{smallmatrix}\right)$, whose mapping torus
on the once-punctured torus is the figure-eight knot complement, which is why
this orbit is known by its braid type. Conjugating the square of that braid by
$\sigma_1$ returns the vineyard word letter for letter, so the two close to the
same knot and the vineyard carries exactly twice the topological entropy,
$1.924847$ against $0.962424$. And unlike the trajectory braid, whose word
depends on the projection chosen, the vineyard braid needs none.

\vspace{0.35em}
{\footnotesize A log behind every number above is at
\texttt{github.com/shawnray-research/natural-monodromy}.}

\end{document}
